\documentclass[10pt,letterpaper]{article}

\usepackage[margin=0.42in]{geometry}
\usepackage{enumitem}
\usepackage[hidelinks]{hyperref}
\usepackage{titlesec}
\usepackage{xcolor}

\setlength{\parindent}{0pt}
\setlength{\parskip}{0pt}
\renewcommand{\baselinestretch}{0.94}
\pagenumbering{gobble}

\titleformat{\section}{\normalsize\bfseries}{}{0pt}{}[\vspace{-0.42em}\rule{\linewidth}{0.35pt}]
\titlespacing*{\section}{0pt}{0.42em}{0.22em}
\setlist[itemize]{leftmargin=1.05em,itemsep=0.04em,topsep=0.03em}
\setlist[enumerate]{leftmargin=1.2em,itemsep=0.06em,topsep=0.03em}

\newcommand{\entry}[2]{\textbf{#1}\hfill #2}
\newcommand{\link}[2]{\href{#1}{#2}}

\begin{document}
\fontsize{9.4}{10.2}\selectfont

\begin{center}
  {\Large \textbf{Kai Ren}}\\[-0.15em]
  Kyoto University \textbar{} Ph.D. Candidate \textbar{} Intelligent Science and Technology\\
  EMG Sensing \textbar{} Robotics \textbar{} Human-Computer Interaction\\
  \href{mailto:ren.kai.34k@st.kyoto-u.ac.jp}{ren.kai.34k@st.kyoto-u.ac.jp}\\
  \href{https://adamrek.github.io}{adamrek.github.io}
\end{center}
\vspace{-0.4em}

\section*{Profile}
Ph.D. candidate in Intelligent Science and Technology at Kyoto University, focusing on electromyography (EMG) sensing, deep learning, robotics, and human-computer interaction. Current work targets physiological signal modeling and motion intent prediction for seat-based and wearable assistive systems, with end-to-end experience from data acquisition to real-time system integration.

\section*{Research Experience}
\entry{RIKEN, Man-Machine Collaboration Research Team, Japan}{Dec 2021 -- Present}\\
\textit{Research Trainee} \hfill \link{https://grp.riken.jp/labs/man_machine_collab/}{RIKEN-GRP Homepage}
\begin{itemize}
  \item Conduct EMG-based analysis and prediction for sit-to-stand and stand-to-sit transitions in assistive scenarios.
  \item Design real-time physiological signal processing pipelines and neural network models for motion intent prediction.
  \item Build experimental closed-loop systems for powered assistive devices, including seat-based prototypes.
  \item Execute the full research workflow: EMG hardware setup, data collection, model training, online evaluation, and validation.
\end{itemize}

\section*{Education}
\begin{itemize}
  \item Kyoto University, Japan, Ph.D. in Intelligent Science and Technology \hfill Apr 2021 -- Present
  \item Ocean University of China, China, M.Eng. in Computer Application Engineering \hfill 2017 -- 2020
  \item Ocean University of China, China, B.Eng. in Computer Science and Technology \hfill 2013 -- 2017
\end{itemize}

\section*{Activities and Projects}
\begin{itemize}
  \item 2014 ACM-ICPC Shandong Provincial Contest participant.
  \item 2016 Project Lead, SRDP Program, Ocean University of China, iOS app development.
  \item 2025 IST Best Poster Award, Graduate School of Informatics, Kyoto University.
  \item Third-Class Academic Scholarship, 2017--2020, awarded consecutively each year.
\end{itemize}

\section*{Publications}
\begin{enumerate}
  \item Anticipatory prediction of sit-to-stand and stand-to-sit transitions: a unified approach. \textit{Frontiers in Bioengineering and Biotechnology}, 2026. DOI: \href{https://doi.org/10.3389/fbioe.2026.1792582}{10.3389/fbioe.2026.1792582}.
  \item Integrated Motion State Prediction for Sit-to-Stand and Stand-to-Sit Motions Toward Effective Power Assist Control. \textit{IEEE ICRA}, 2025, pp. 947--952.
  \item Q-DFN: Q-learning Decisive Fusion Network for Object Detection. \textit{Proceedings of CSAI 2019}, pp. 88--92.
\end{enumerate}

\section*{Technical Skills}
\begin{itemize}
  \item \textbf{Programming:} Python, MATLAB, C++
  \item \textbf{Machine Learning:} PyTorch, deep learning for biosignals and time-series modeling
  \item \textbf{Signal Processing:} EMG acquisition, preprocessing, feature extraction, real-time processing
  \item \textbf{Systems and Hardware:} real-time inference, closed-loop systems, assistive-device integration, EMG sensor setup
  \item \textbf{Tools:} Linux, Git, experimental system integration
\end{itemize}

\section*{Research Interests and Languages}
\textbf{Research Interests:} physiological signal modeling for motion intent prediction, multimodal perception, real-time reasoning, and human-centered AI.\\
\textbf{Languages:} Chinese, native; English, TOEFL; Japanese, JLPT N2.

\end{document}
